% !TEX program = pdflatex
\documentclass[11pt]{article}

\usepackage[margin=1in]{geometry}
\usepackage{amsmath,amssymb}
\usepackage{booktabs}
\usepackage{longtable}
\usepackage{array}
\usepackage{hyperref}
\usepackage{xurl}
\usepackage{microtype}
\usepackage{enumitem}

\title{Natural and Vaccine-Induced Immunity:\\
Approximate Waning Rates for Major Vaccine-Preventable Infections}
\author{}
\date{September 2026 (revised)}

\begin{document}
\maketitle

\begin{abstract}
This note compares the duration of protection following natural infection with
the duration of protection following vaccination for a selection of important
infectious diseases, including COVID-19. The information is based primarily on
official material from the U.S. Centers for Disease Control and Prevention
(CDC) --- chiefly the \emph{Pink Book}, CDC's reference text on vaccine-preventable
diseases --- together with World Health Organization (WHO) material where
appropriate.

A numerical ``mean duration of immunity'' is not available for many diseases.
Official sources frequently use expressions such as ``lifelong,'' ``probably
lifelong,'' ``at least 20 years,'' or ``waning over several years.'' It would
therefore be misleading to assign a spurious numerical mean to all diseases.
The first part of the note reports the official qualitative or observational
estimates, disease by disease. In the final table, these are converted --
where, and only where, this is reasonably justified -- into an indicative
waning rate
\[
        \lambda = \frac{1}{T},
\]
measured in $\mathrm{day}^{-1}$, where $T$ is a characteristic duration in
days. Every table entry is coded using one of four symbols, defined
immediately below the table: \textbf{0} (effectively lifelong, hence no
finite waning rate), a \textbf{numerical value}, \textbf{N/D} (not
determined precisely enough by the evidence), or \textbf{N/C} (not
meaningfully reducible to a single characteristic rate).
\end{abstract}

\section{Interpretation of the comparison}

It is important to distinguish three related but different quantities:

\begin{enumerate}
\item protection against infection or reinfection;
\item protection against symptomatic disease; and
\item protection against severe disease, hospitalization, or death.
\end{enumerate}

These quantities can have quite different durations. This is particularly
important for influenza and COVID-19, for which immunity against infection
can wane considerably faster than protection against severe disease.

Moreover, ``waning'' is not generally a simple exponential process. The
quantity $\lambda = 1/T$ used in the table below is therefore a
\emph{characteristic inverse timescale}, not a directly measured biological
decay constant. A genuine exponential model would require a well-defined
survival or protection curve and would normally use a hazard or decay
parameter estimated from longitudinal data. For several diseases discussed
below, official sources decline to provide any such curve, and in those
cases no number is invented.

\section{Summary of official evidence}

\subsection{Measles}

The CDC states that immunity following natural measles infection is
lifelong; this observation dates at least to Peter Panum's 1846 description
of the disease (Ref.~\ref{ref:pinkbook-ch13}). For measles-containing
vaccine, the CDC states that both serologic and epidemiologic evidence
indicate that vaccine-induced immunity is long-term and probably lifelong in
most persons. After two doses, more than 99\% of persons develop serologic
evidence of immunity, and persons who appear to lose antibody typically show
a rapid (anamnestic) response on re-exposure, suggesting that protection
persists even when circulating antibody is low (Ref.~\ref{ref:pinkbook-ch13}).

Thus, for both natural infection and vaccination, the appropriate
description is essentially lifetime protection rather than a finite mean
duration.

\subsection{Mumps}

Mumps is the clearest exception among the MMR-preventable diseases. Natural
mumps infection is generally described as conferring lifelong immunity, but
this characterization is somewhat softer than for measles or rubella:
reinfections, though uncommon, are recognized to occur.

Vaccine-induced immunity is where mumps differs most sharply from measles
and rubella. The CDC's Pink Book chapter on mumps quantifies vaccine
\emph{effectiveness} (78\% after one dose, 88\% after two doses) but --
unlike the measles, rubella, or polio chapters -- does not offer any
official statement of the \emph{duration} of protection
(Ref.~\ref{ref:pinkbook-ch15}). CDC's current public-facing summary states
only that ``immunity against mumps may decrease over time''
(Ref.~\ref{ref:cdc-measles-vax}). In the absence of an official duration, we
draw on a widely cited peer-reviewed academic estimate: pooling six
vaccine-effectiveness studies with a dynamic transmission model, Lewnard and
Grad estimated that vaccine-derived protection against mumps wanes with a
mean duration of 27 years (95\% confidence interval, 16--51 years) after
vaccination, and that this waning -- rather than viral evolution -- explains
the resurgence of mumps among vaccinated young adults in the United States
since 2006 (Ref.~\ref{ref:lewnard-grad}). This is an academic modeling
estimate, not an official CDC or WHO figure, and it is flagged as such in
the table below.

\subsection{Rubella}

The CDC reports that more than 90\% of vaccinated persons have protection
against clinical rubella for at least 15 years, and that follow-up studies
indicate that one dose provides long-term, probably lifelong, protection.
Although antibody titers wane in the years after vaccination, the CDC notes
there is no evidence this leads to significant susceptibility to clinical
rubella or congenital rubella syndrome (Ref.~\ref{ref:pinkbook-ch20}).
Natural infection is generally regarded as producing long-lasting immunity.

Because the CDC's own bottom-line characterization of the vaccine is
``probably lifelong'' -- with the 15-year figure serving as supporting
evidence for that conclusion rather than as the conclusion itself -- rubella
vaccine, like the measles vaccine, is treated below as falling in the
``effectively lifelong'' category rather than being assigned a finite rate.
This differs from the earlier draft of this note, which had, inconsistently,
converted the 15-year figure directly into a rate while treating comparably
worded statements for other diseases as indicating zero waning; the revision
here applies one rule uniformly, and it is discussed further below the
table.

\subsection{Varicella}

Following vaccination, the CDC reports that immunity appears long-lasting
and is probably permanent in the majority of recipients. More than 90\% of
vaccine responders maintain antibody for at least six years, with antibody
also documented seven to ten years after vaccination in Japanese studies
(Ref.~\ref{ref:pinkbook-ch22}). As with rubella, these interim data points
support the CDC's ``probably permanent'' conclusion rather than replacing it
with a finite number.

Natural varicella infection generally produces lifetime immunity, although
the precise duration is not represented by a single experimentally measured
mean; second occurrences of chickenpox are uncommon in immunocompetent
persons.

\subsection{Hepatitis A}

The CDC states plainly that the exact duration of protection following
hepatitis A vaccination is unknown -- a materially different, more
open-ended statement than the ``probably lifelong'' language used for
measles, rubella, varicella, or polio. Anti-HAV has persisted for at least
25 years in adults vaccinated as children under the older three-dose
schedule used before 1999, and for at least 20 years in adults who received
a two-dose schedule. Mathematical modeling and antibody-kinetic studies
estimate that detectable antibody could persist for 40 years or longer
(Ref.~\ref{ref:pinkbook-ch9}). Because the CDC explicitly withholds a
``lifelong'' characterization here, these figures are used below as the
basis for a genuine (upper-bound) numerical rate.

Natural hepatitis A infection is generally regarded as producing durable,
essentially lifelong immunity.

\subsection{Hepatitis B}

Hepatitis B vaccination produces long-lasting immunologic memory. The CDC
reports that immune memory remains intact for more than 30 years following
immunization, and that both adults and children with declining antibody
levels remain protected against clinically significant infection because of
a rapid anamnestic response on exposure. Booster doses are not recommended
for immunocompetent persons (Ref.~\ref{ref:pinkbook-ch10}).

The interpretation of \emph{natural} immunity is fundamentally more
complicated, because infection can either resolve with immunity or become
chronic, and these are qualitatively different outcomes rather than two
points on the same waning curve. For this reason -- discussed further below
-- natural hepatitis B immunity is not assigned a single characteristic
rate at all; the duration of protective immunity after a resolved infection
should in any case not be confused with the natural history of chronic
hepatitis B infection itself.

\subsection{Poliomyelitis}

Natural poliovirus infection produces long-lasting immunity. The CDC
describes immunity after a complete vaccination series -- with either
inactivated (IPV) or live oral (OPV) vaccine -- as probably lifelong,
while noting that the precise duration cannot be experimentally pinned down
to a single mean value (Ref.~\ref{ref:pinkbook-ch18}).

\subsection{Pertussis}

Pertussis provides a particularly useful contrast. The CDC states
explicitly that immunity following \emph{Bordetella pertussis} infection is
not permanent (Ref.~\ref{ref:pinkbook-ch16}). Published epidemiological
estimates place the duration of natural protection in the range of roughly
four to twenty years, depending on the study and the definition of
protection used.

Vaccine-induced immunity is shorter still. Protection following
acellular pertussis vaccination (DTaP/Tdap) wanes appreciably over
approximately four to five years, which is a principal reason booster doses
are recommended (Ref.~\ref{ref:pinkbook-ch16}).

\subsection{Influenza}

Influenza differs fundamentally from the diseases above because immunity is
shaped not only by waning immune responses but also by the antigenic
evolution (``drift'') of the circulating virus. The CDC states explicitly,
as an official practical conclusion, that ``for practical purposes, the
duration of immunity following influenza vaccination is less than one year
because of vaccine-induced antibody waning and antigenic drift of
circulating influenza viruses'' (Ref.~\ref{ref:pinkbook-ch12}). Because this
is an explicit, official, single-number statement, it is used below as a
genuine (if composite) rate for the vaccine.

No comparable official statement exists for the duration of \emph{natural}
post-infection immunity to influenza. This is not merely an oversight: a
single number would conflate true immunological waning (which, against an
unchanged strain, may be prolonged) with antigenic drift of the virus
itself (which can defeat even fully intact immunity within a season). In
the absence of an official figure isolating the two effects, natural
influenza immunity is left undetermined below rather than assumed equal to
the vaccine figure, as the earlier draft of this note had done.

\subsection{Human papillomavirus (HPV)}

Natural HPV infection does not reliably produce durable protective immunity
against subsequent infection; reinfection and infection with additional
HPV types can occur. This is a qualitatively different situation from
``fast waning'' -- for many infections, robust cross-protective immunity
is never clearly established in the first place, so there is no single
protection curve to characterize.

HPV vaccination, by contrast, produces strong and durable antibody
responses. The CDC's Pink Book states that ongoing monitoring has
demonstrated vaccine effectiveness above 90\%, with no waning of immunity
through at least 10 to 12 years of follow-up (Ref.~\ref{ref:pinkbook-ch11}),
consistent with CDC's public-facing statement that protection lasts more
than 10 years without becoming less effective and that no data indicate
protection lessens over time (Ref.~\ref{ref:cdc-hpv-safety}). As with
hepatitis A, this is phrased as an absence of observed waning within a
follow-up window, not as an affirmative ``lifelong'' finding, so it is used
below as the basis for a genuine (upper-bound) numerical rate rather than
being coded as zero.

\subsection{COVID-19}

COVID-19 requires particular caution. Natural infection produces immune
responses in most individuals, but protection against reinfection changes
with time and with the emergence of antigenically different variants. The
WHO's 2021 scientific brief on natural immunity reported that, in the
evidence available at the time, most infected people had robust protective
immune responses for at least 6--8 months, while explicitly cautioning that
this finding ``should not be interpreted as a fixed biological expiration
date'' and flagging that variants of concern could alter the picture
(Ref.~\ref{ref:who-brief}).

Vaccine-induced protection shows a similar pattern but is explicitly
bifurcated by endpoint in official communications: protection against
infection wanes over a period of months, while protection against severe
disease, hospitalization, and death generally persists substantially
longer. Updated vaccines are recommended precisely because both immune
waning and antigenic evolution alter protection over time.

Unlike influenza, no official source commits to a single practical duration
figure for COVID-19 that spans both endpoints and both exposure routes.
For this reason, neither natural nor vaccine-induced COVID-19 immunity is
reduced to a single rate in the table below; both are coded as not
meaningfully comparable via one characteristic waning rate, which is a
stronger conclusion than simply lacking a number, and is consistent with
how official sources themselves discuss COVID-19 immunity.

\subsection{A note on the stability of ``official'' sources in 2025--2026}

Between 2025 and 2026, U.S. federal vaccine recommendations went through an
unusually turbulent period: the Secretary of Health and Human Services
replaced the full membership of the Advisory Committee on Immunization
Practices (ACIP) in 2025, ACIP subsequently revised several recommendations
(including moving the birth-dose hepatitis B recommendation from universal
to shared clinical decision-making, discontinuing preference for the
combination MMRV vaccine, and recommending a single-dose HPV schedule), and
a federal court stayed both the resulting 2026 childhood immunization
schedule and the disputed ACIP appointments pending litigation
(Ref.~\ref{ref:crs-acip}). These disputes concern vaccination
\emph{schedules and policy}, not the underlying serologic and
epidemiologic evidence on duration of protection summarized above, which
predates this period and is unaffected by it. Readers should nonetheless be
aware that ``the official CDC position'' has been an unusually moving
target in this specific window, and should check current primary sources
for anything schedule-related.

\section{Why inverse duration is used}

Suppose, purely as a phenomenological approximation, that a protection
timescale is represented by $T$ days. We define the corresponding
characteristic waning rate by
\[
       \boxed{\lambda = \frac{1}{T}\quad \mathrm{day}^{-1}}.
\]

This definition is useful for comparing orders of magnitude. For example,

\[
\begin{array}{c|c}
T & 1/T\;(\mathrm{day}^{-1})\\
\hline
1\ {\rm year\ (365\ days)} & 2.74\times 10^{-3}\\
5\ {\rm years} & 5.48\times 10^{-4}\\
10\ {\rm years} & 2.74\times 10^{-4}\\
20\ {\rm years} & 1.37\times 10^{-4}\\
27\ {\rm years} & 1.01\times 10^{-4}\\
40\ {\rm years} & 6.85\times 10^{-5}
\end{array}
\]

(All year-to-day conversions in this note use 365 days per year, for
consistency; this introduces no meaningful error at the precision used
here.) These numbers should be read as inverse characteristic durations,
not as experimentally established exponential decay constants.

For lifelong immunity, it is more scientifically honest to write ``the
waning rate is effectively zero on the finite epidemiological timescales
considered'' than to assign an arbitrary human life expectancy as $T$.
Accordingly, the numerical table below uses the symbol $0$ whenever the
official source's own bottom-line characterization is ``lifelong'' or
``probably lifelong,'' even if a finite interim duration is cited elsewhere
in the source as supporting evidence for that conclusion.

\section{Numerical comparison}

The table below is deliberately conservative and follows one rule
throughout: a number is entered only when the cited source's own
substantive conclusion is a finite, bounded duration -- not when a finite
duration merely appears as an intermediate data point supporting a
``lifelong'' conclusion, and not when the underlying process cannot be
sensibly reduced to a single rate at all. Every cell in the last three
columns is coded using exactly one of four symbols, defined immediately
below the table. Diseases are grouped by the value of the ratio in the
final column, in the order $0$, then a nonzero number, then N/C, then N/D.

\begin{longtable}{>{\raggedright\arraybackslash}p{2.0cm}
                  >{\raggedright\arraybackslash}p{2.8cm}
                  >{\raggedright\arraybackslash}p{2.75cm}
                  >{\raggedleft\arraybackslash}p{1.9cm}
                  >{\raggedleft\arraybackslash}p{1.9cm}
                  >{\raggedleft\arraybackslash}p{1.7cm}}
\caption{Indicative waning rates, writing $\theta \equiv \lambda_{\rm
natural} = 1/T_{\rm natural}$ for the characteristic natural-immunity
waning rate and $\psi \equiv \lambda_{\rm vaccine} = 1/T_{\rm vaccine}$ for
the characteristic vaccine-immunity waning rate (both in
$\mathrm{day}^{-1}$). Diseases are grouped by the value of $\theta/\psi$,
coded as described below the table.}
\label{tab:waning}\\
\toprule
Disease & Natural immunity: official characterization & Vaccine immunity:
official characterization & Natural $\theta$ & Vaccine $\psi$
& $\theta/\psi$\\
\midrule
\endfirsthead
\toprule
Disease & Natural immunity: official characterization & Vaccine immunity:
official characterization & Natural $\theta$ & Vaccine $\psi$
& $\theta/\psi$\\
\midrule
\endhead
\midrule \multicolumn{6}{r}{\emph{(continued on next page)}}\\ \midrule
\endfoot
\bottomrule
\endlastfoot

\multicolumn{6}{l}{\textbf{Group 1: $\theta/\psi = 0$} (natural immunity effectively non-waning; vaccine wanes at a finite rate)}\\
\midrule

Mumps &
Generally lifelong, but reinfection documented more often than for measles/rubella &
No official duration; academic model: mean 27 yr (95\% CI 16--51 yr) &
$0$ &
$1.0\times10^{-4}$ \newline (0.5--1.7$\times10^{-4}$) &
$0$ \\

Hepatitis A &
Essentially lifelong &
CDC: exact duration unknown; documented $\geq$20--25 yr; modeled $\geq$40 yr &
$0$ &
$\leq 1.37\times10^{-4}$ \newline (as low as $6.85\times10^{-5}$ under modeling) &
$0$\\

\midrule
\multicolumn{6}{l}{\textbf{Group 2: $\theta/\psi$ a nonzero number} (both mechanisms wane at a finite rate)}\\
\midrule

Pertussis &
CDC: not permanent; published epidemiological estimates $\approx$4--20 yr &
Wanes appreciably over $\approx$4--5 yr (acellular vaccine) &
$1.37\times10^{-4}$--$6.85\times10^{-4}$ &
$5.48\times10^{-4}$--$6.85\times10^{-4}$ &
$\approx0.2$--$1.25$ \\

\midrule
\multicolumn{6}{l}{\textbf{Group 3: $\theta/\psi =$ N/C} (not meaningfully comparable via one characteristic rate)}\\
\midrule

Measles &
Lifelong (Panum, 1846) &
Long-term, probably lifelong; $>$99\% seroconvert after 2 doses &
$0$ &
$0$ &
N/C \\

Rubella &
Long-lasting / essentially lifelong &
$\geq$15 yr documented; officially ``probably lifelong'' &
$0$ &
$0$ &
N/C \\

Varicella &
Long-lasting; second episodes uncommon &
Antibody documented $\geq$6--10 yr; officially ``probably permanent'' &
$0$ &
$0$ &
N/C \\

Poliomyelitis &
Long-lasting &
IPV and OPV both officially ``probably lifelong'' &
$0$ &
$0$ &
N/C\\

Hepatitis B &
Heterogeneous: resolves with immunity, or becomes chronic; not one process &
Immune memory documented $\geq$30 yr; boosters not advised in immunocompetent adults &
N/C &
$\leq 9.13\times10^{-5}$ &
N/C \\

HPV &
No reliable durable natural protection; not a simple waning process &
CDC: $>$90\% effective, no waning observed through 10--12 yr &
N/C &
$\leq2.74\times10^{-4}$ &
N/C \\

COVID-19 &
WHO (2021): most had robust responses $\geq$6--8 mo against reinfection; explicitly not a fixed expiry; variant-dependent &
Protection vs.\ infection wanes over months; vs.\ severe disease persists longer; no single official figure &
N/C &
N/C &
N/C \\

\midrule
\multicolumn{6}{l}{\textbf{Group 4: $\theta/\psi =$ N/D} (not determined precisely enough by official evidence)}\\
\midrule

Influenza &
No official duration; confounded with antigenic drift &
CDC: practically $<$1 yr, from waning \emph{and} antigenic drift &
N/D &
$\gtrsim2.74\times10^{-3}$ &
N/D \\

\end{longtable}

\subsection*{Key to the table}

\begin{description}[leftmargin=2.2cm,style=nextline]
\item[$0$] Effectively lifelong protection, hence a zero waning rate on the
epidemiological timescales considered here. Used whenever the cited
source's own substantive conclusion is ``lifelong'' or ``probably
lifelong,'' even if a shorter documented interval is cited elsewhere in the
same source as supporting evidence.
\item[numerical value] A defensible finite characteristic rate. Used only
when the cited source's own conclusion is itself a bounded, finite
duration (an observed floor, a modeled estimate, or an explicit practical
figure) rather than an affirmation of durability. Where the underlying
duration is itself a range or a confidence interval, the corresponding
range of $\theta$ or $\psi$ is given.
\item[N/D] Not determined sufficiently well by the official evidence. Used
when a finite characteristic duration plausibly exists in principle but no
official source states or bounds it (natural influenza immunity is the
clearest case here).
\item[N/C] Not meaningfully comparable using one characteristic waning
rate. Used when the underlying phenomenon cannot be sensibly reduced to a
single decay parameter at all -- for example, because natural infection
outcomes are qualitatively heterogeneous (hepatitis B), because natural
infection does not reliably establish protective immunity in the first
place (HPV), because two protection endpoints diverge too far to summarize
in one number (COVID-19), or because a ratio would otherwise involve
dividing by, or into, an undefined or ill-posed quantity.
\end{description}

A methodological note: several of these classifications required a
judgment call, and a reasonable person could draw some of these lines
differently -- most notably for rubella and varicella vaccine, discussed in
Sections 2.3 and 2.4 above, where an interim documented duration exists
alongside an official ``probably lifelong'' conclusion. The rule applied
throughout this revision is to code according to the source's own
substantive conclusion, not the most conservative number appearing
anywhere in its supporting text; this is a change from the earlier version
of this note, which applied that rule inconsistently across diseases.

The ratio column, $\theta/\psi$, follows directly from the coding of its
numerator $\theta$ and denominator $\psi$: a ratio of exactly $0$ is entered
only when $\theta$ is coded $0$ and $\psi$ is a genuine finite number
(giving a mathematically exact zero, not merely an approximation); the
ratio is coded N/C whenever \emph{either} $\theta$ or $\psi$ is itself $0$
paired with $0$, or is coded N/C; and it is coded N/D when either component
is coded N/D. Note that this differs from, and corrects, the earlier
version of this table, which entered ``$\infty$'' for several diseases
where both rates were effectively zero -- a case where $\theta/\psi$ is
actually indeterminate ($0/0$), not infinite.

\paragraph{Important qualification.} The COVID-19 and influenza entries
require special care. Even where a number is given for influenza vaccine
(the one case with an explicit official practical figure), it reflects a
mixture of true immunological waning and antigenic drift of the circulating
virus, not a pure decay constant. The COVID-19 figures discussed in the
text above are deliberately \emph{not} carried into the table as numbers,
precisely because official sources themselves decline to summarize
COVID-19 immunity with one figure; none of this should be read as implying
that protection disappears abruptly at any of the quoted intervals.

\section{Official references}

\begin{enumerate}

\item \label{ref:pinkbook-ch1} CDC, \emph{Pink Book, Chapter 1: Principles
of Vaccination}. Explains active immunity, immunologic memory, and the
general distinction between immunity acquired through infection and
through vaccination.
\url{https://www.cdc.gov/pinkbook/hcp/table-of-contents/chapter-1-principles-of-vaccination.html}

\item \label{ref:pinkbook-ch9} CDC, \emph{Pink Book, Chapter 9: Hepatitis
A}. States that the exact duration of protection after vaccination is
unknown, and reports persistence of anti-HAV for at least 20--25 years and
model-based estimates of 40 years or longer.
\url{https://www.cdc.gov/pinkbook/hcp/table-of-contents/chapter-9-hepatitis-a.html}

\item \label{ref:pinkbook-ch10} CDC, \emph{Pink Book, Chapter 10: Hepatitis
B}. Reports that immune memory remains intact for more than 30 years after
vaccination and that booster doses are not recommended for immunocompetent
persons.
\url{https://www.cdc.gov/pinkbook/hcp/table-of-contents/chapter-10-hepatitis-b.html}

\item \label{ref:pinkbook-ch11} CDC, \emph{Pink Book, Chapter 11: Human
Papillomavirus}. Reports vaccine effectiveness above 90\% with no waning of
immunity observed through at least 10 to 12 years of follow-up.
\url{https://www.cdc.gov/pinkbook/hcp/table-of-contents/chapter-11-human-papillomavirus.html}

\item \label{ref:cdc-hpv-safety} CDC, \emph{HPV Vaccine Safety and
Effectiveness Data}. States that protection lasts more than 10 years
without becoming less effective and that no data indicate waning over time.
\url{https://www.cdc.gov/hpv/hcp/vaccination-considerations/safety-and-effectiveness-data.html}

\item \label{ref:pinkbook-ch12} CDC, \emph{Pink Book, Chapter 12:
Influenza}. States that, for practical purposes, the duration of immunity
following influenza vaccination is less than one year, due to a
combination of antibody waning and antigenic drift.
\url{https://www.cdc.gov/pinkbook/hcp/table-of-contents/chapter-12-influenza.html}

\item \label{ref:pinkbook-ch13} CDC, \emph{Pink Book, Chapter 13:
Measles}. States that vaccine-induced measles immunity is long-term and
probably lifelong in most persons, and that natural infection confers
lifelong immunity.
\url{https://www.cdc.gov/pinkbook/hcp/table-of-contents/chapter-13-measles.html}

\item \label{ref:pinkbook-ch15} CDC, \emph{Pink Book, Chapter 15: Mumps}.
Reports vaccine effectiveness (78\% after one dose, 88\% after two doses)
but does not state an official duration of protection.
\url{https://www.cdc.gov/pinkbook/hcp/table-of-contents/chapter-15-mumps.html}

\item \label{ref:pinkbook-ch16} CDC, \emph{Pink Book, Chapter 16:
Pertussis}. States explicitly that immunity after \emph{B.\ pertussis}
infection is not permanent, and discusses waning of acellular vaccine
immunity over several years.
\url{https://www.cdc.gov/pinkbook/hcp/table-of-contents/chapter-16-pertussis.html}

\item \label{ref:pinkbook-ch18} CDC, \emph{Pink Book, Chapter 18:
Poliomyelitis}. Describes immunity from a complete IPV or OPV series as
probably lifelong.
\url{https://www.cdc.gov/pinkbook/hcp/table-of-contents/chapter-18-poliomyelitis.html}

\item \label{ref:pinkbook-ch20} CDC, \emph{Pink Book, Chapter 20:
Rubella}. Reports protection for at least 15 years and follow-up
suggesting one dose confers long-term, probably lifelong, protection.
\url{https://www.cdc.gov/pinkbook/hcp/table-of-contents/chapter-20-rubella.html}

\item \label{ref:pinkbook-ch22} CDC, \emph{Pink Book, Chapter 22:
Varicella}. Describes vaccine-induced immunity as long-lasting and
probably permanent in the majority of recipients.
\url{https://www.cdc.gov/pinkbook/hcp/table-of-contents/chapter-22-varicella.html}

\item \label{ref:pinkbook-index} CDC, \emph{Pink Book, Site Index}. Full
table of contents for the Pink Book's chapters on individual
vaccine-preventable diseases.
\url{https://www.cdc.gov/pinkbook/site.html}

\item \label{ref:cdc-measles-vax} CDC, \emph{Measles Vaccination}.
States that MMR/MMRV vaccines usually protect against measles and rubella
for life, while immunity against mumps may decrease over time.
\url{https://www.cdc.gov/measles/vaccines/index.html}

\item \label{ref:who-brief} WHO, \emph{COVID-19 Natural Immunity:
Scientific Brief}, 10 May 2021. Reports that most infected persons had
robust protective immune responses for at least 6--8 months in the
evidence available at that time, while cautioning against interpreting
this as a fixed expiration date.
\url{https://iris.who.int/handle/10665/341241}

\item \label{ref:lewnard-grad} Lewnard JA, Grad YH. Vaccine waning and
mumps re-emergence in the United States. \emph{Sci Transl Med}.
2018;10(433):eaao5945. Peer-reviewed academic estimate (not an official
CDC/WHO figure) that vaccine-derived mumps immunity wanes with a mean
duration of 27 years (95\% CI, 16--51 years).
\url{https://doi.org/10.1126/scitranslmed.aao5945}

\item \label{ref:crs-acip} Congressional Research Service, \emph{Changes
to CDC Vaccine Recommendations in 2025 and 2026}. Summarizes ACIP
membership changes, the resulting revisions to vaccine recommendations,
and related litigation over the 2026 childhood immunization schedule.
\url{https://www.congress.gov/crs-product/IN12684}

\end{enumerate}

\section{Concluding observation}

The principal lesson from the comparison is that the ratio $\theta/\psi$
(writing $\theta \equiv \lambda_{\rm natural}$ and $\psi \equiv
\lambda_{\rm vaccine}$, as in Table~\ref{tab:waning}), where it can be
formed at all, is highly disease-dependent. There is no universal
biological principle according to which natural immunity must wane more
slowly than vaccine-induced immunity, nor the converse.

For measles, rubella, varicella, and polio, both natural infection and
vaccination are officially characterized as producing extremely durable,
essentially non-waning immunity, and the ratio between them is not
meaningfully defined. For mumps and hepatitis A, natural infection is
essentially non-waning while the vaccine wanes on a finite (multi-decade)
timescale, giving a ratio of exactly zero. For pertussis, both mechanisms
wane on a finite timescale, with natural immunity lasting somewhat longer
on average than protection from the current acellular vaccines, though
neither is permanent. For HPV, natural infection is a poor and inconsistent
source of durable protection, whereas vaccination produces very durable
protection -- the opposite pattern from pertussis, and one where the
natural-immunity side of the ratio is not well-posed at all. For hepatitis
B, the very idea of a single natural-immunity waning rate breaks down
because infection outcomes are qualitatively heterogeneous. For influenza
and COVID-19, waning immunity and antigenic evolution are entangled to the
point that official sources themselves decline to offer a single duration,
and this note follows their lead rather than manufacturing false precision.

Consequently, the inverse-duration table is best regarded as a compact
phenomenological summary useful for modeling in the cases where the
underlying evidence supports it, and as an explicit record of where it does
not, rather than as a table of fundamental immunological decay constants.

\end{document}